% =====================================================================
% Capítulo 1 - Introducción
% Tesis: Extensión de ST-CC con refinamiento jerárquico por bloques
%        para la detección de patrones espacio-temporales locales
% =====================================================================

\chapter{Introducción}
\label{chap:introduccion}

El presente capítulo introduce el problema abordado en esta tesis. Se parte
del contexto general del análisis de datos espacio-temporales, se describe
la limitación específica del co-agrupamiento de rejilla regular que motiva
el trabajo, y se presentan los objetivos, la organización del documento y
el cronograma de actividades.

% ---------------------------------------------------------------------
\section{Motivación y Contexto}
\label{sec:motivacion}

El avance de los sensores, los dispositivos con posicionamiento global y
los sistemas de registro automático ha hecho que grandes volúmenes de datos
espacio-temporales estén disponibles de forma cotidiana. Los registros de
viajes en taxi, los catálogos sísmicos o las series de monitoreo ambiental
describen dónde y cuándo ocurre un fenómeno, y su análisis
permite comprender la dinámica de una ciudad o de un territorio y apoyar la
toma de decisiones en planificación urbana, transporte y gestión de riesgos
\cite{Liu2021}.

Sin embargo, estos datos presentan una dificultad característica, y es que la
dimensión espacial y la temporal no son independientes. En la movilidad
urbana, por ejemplo, los destinos de los viajes se concentran en zonas de
trabajo durante la mañana y en zonas residenciales durante la tarde. Un
patrón de este tipo solo es visible si se observan ambas dimensiones a la
vez. Las técnicas tradicionales de agrupamiento espacio-temporal suelen
analizar una sola dimensión, o bien tratan cada dimensión por separado, lo
cual puede producir resultados poco representativos y perder la información
de interacción entre espacio y tiempo \cite{Liu2021}.

El \ac{CC} ofrece una alternativa natural a este problema. En lugar de
agrupar solo filas o solo columnas, particiona simultáneamente ambas
dimensiones de una matriz de datos. Aplicado a una matriz de co-ocurrencia
$O \in \mathbb{R}_{+}^{m \times n}$, cuyas filas son regiones y cuyas
columnas son intervalos de tiempo, el co-agrupamiento identifica grupos de
regiones que se comportan de manera similar en ciertos grupos de horarios.
Sobre esta idea, \cite{Liu2021} propusieron el marco \ac{ST-CC}, que emplea
el algoritmo \ac{BBACI} \cite{Banerjee2007} para descubrir patrones de
movilidad en los datos de taxis de Manhattan. Su estudio utilizó alrededor
de diez millones de viajes registrados en noviembre de 2017, agregados en
66 regiones y 12 intervalos de dos horas.

El resultado de \ac{BBACI} es una partición en tablero de ajedrez
(\emph{checkerboard}) donde cada región pertenece a un único grupo de regiones y
cada intervalo a un único grupo temporal, de modo que la matriz queda
dividida en $K_R \times K_T$ bloques rectangulares. Esta estructura es
fácil de interpretar, pero es global. Si una subzona de un grupo de
regiones se comporta de manera distinta solo en un horario concreto, por
ejemplo un foco de actividad alrededor de una estación de trenes al final
de la tarde, la única forma de capturarlo es dividir el grupo de regiones
completo en todos los horarios. En otras palabras, un detalle local obliga
a aumentar la resolución de toda la fila o columna de la partición. Este
problema también fue señalado en la literatura de modelos de bloques, donde
motivó enfoques bayesianos no paramétricos como el proceso de Mondrian
\cite{Roy2008}.

\cite{Liu2021} reconocieron parte de esta limitación y aplicaron
K-means sobre los bloques del tablero para obtener \emph{co-clusters} que no
siguen la rejilla. No obstante, ese paso reagrupa bloques ya
existentes a partir de su media y desviación estándar; no subdivide
un bloque internamente. Por tanto, la heterogeneidad que queda oculta dentro
de un bloque, que es precisamente donde pueden residir los patrones locales,
no se recupera. La alternativa de aumentar $K_R$ y $K_T$, o de elegirlos
automáticamente con métodos como Cross-Associations \cite{Chakrabarti2004},
sigue produciendo una rejilla global y aumenta el número de parámetros en
toda la matriz.

Una segunda pregunta surge al intentar refinar localmente, ¿cuándo debería detenerse?
Todo bloque presenta algo de variación, de modo que se necesita un criterio
que distinga la heterogeneidad real del ruido. Dos rasgos de los datos de
movilidad hacen que esa distinción no sea trivial. El primero es que la matriz
agrega varios días, y la actividad de cada celda varía de un día a otro mucho
más de lo que supone el modelo de Poisson, y de forma correlacionada entre
celdas, porque factores como el clima afectan a muchas regiones a la vez; una
prueba que ignore esa variabilidad declara heterogéneo casi cualquier bloque. El
segundo es que, con millones de viajes, incluso una diferencia irrelevante
resulta estadísticamente significativa. En agrupamiento interpretable
basado en árboles, \cite{Hu2025} propusieron evaluar cada división candidata
mediante pruebas de hipótesis, elegir la de menor $p$-valor y detener el
crecimiento cuando ninguna división es estadísticamente significativa. Este
principio, sin embargo, fue desarrollado para agrupar datos categóricos en
una sola dimensión y no para bloques bidimensionales de una matriz
espacio-temporal.

Esta tesis se sitúa en la intersección de ambas líneas. Se propone extender
\ac{ST-CC} con un refinamiento jerárquico que actúe dentro de cada
bloque, guiado por un criterio estadístico coherente con la misma función
objetivo que optimiza \ac{BBACI}. En el caso de conteos, la I-divergencia
que minimiza \ac{BBACI} coincide, salvo un factor constante, con la
devianza de un modelo log-lineal de Poisson. En consecuencia, la reducción
de divergencia al partir un bloque puede interpretarse como un estadístico
de razón de verosimilitud. Esta conexión permite construir el criterio de
refinamiento a partir del propio marco de \ac{ST-CC}, en lugar de importar
un criterio ajeno. Para que el criterio sea fiable, su distribución de
referencia se obtiene remuestreando los días que forman la matriz, de modo que
refleje la variabilidad real de los datos y no la que supondría un modelo de
Poisson, y la decisión de refinar exige, además de significancia, un tamaño de
efecto mínimo. El resultado puede aplicarse en otros dominios con datos de
conteo espacio-temporales agregados a partir de réplicas, como días o semanas.

% ---------------------------------------------------------------------
\section{Planteamiento del Problema}
\label{sec:problema}

El marco \ac{ST-CC} impone una partición en rejilla regular sobre la matriz
región~$\times$~tiempo, en la que cada división de filas o columnas se
aplica a la matriz completa. Como consecuencia, los patrones
espacio-temporales locales, es decir, la interacción entre un
subconjunto de regiones y un subconjunto de horarios que ocurre solo dentro
de un bloque, quedan promediados y ocultos, y la única manera de revelarlos
es aumentar la resolución global a costa de fragmentar innecesariamente el
resto de la matriz. Además, no existe un criterio estadístico que indique
si un bloque contiene heterogeneidad real que justifique subdividirlo, ni
que controle el riesgo de producir refinamientos espurios cuando se evalúan
muchas divisiones candidatas sobre los mismos datos. Ese criterio debe tener
en cuenta que la matriz es la suma de varios días cuya variabilidad excede la
de un modelo de Poisson, y que con millones de viajes la significancia
estadística no basta para distinguir un patrón relevante. El problema que aborda
esta tesis es, por tanto, cómo refinar localmente y de forma
estadísticamente justificada la partición de \ac{ST-CC} para detectar
patrones espacio-temporales locales, manteniendo la interpretabilidad del
resultado.

% ---------------------------------------------------------------------
\section{Objetivos}
\label{sec:objetivos}

El objetivo general de esta tesis es proponer y validar una extensión de
\ac{ST-CC} basada en un refinamiento jerárquico por bloques, que se aplique
sobre la partición en tablero que produce el marco original y que esté guiado
por pruebas de significancia estadística derivadas de la I-divergencia. Se
busca que esta extensión permita detectar patrones espacio-temporales locales,
es decir, la interacción entre un subconjunto de regiones y un subconjunto de
horarios que ocurre solo dentro de un bloque y que la partición en rejilla
regular no captura. Al mismo tiempo, se espera que el criterio con el que se
decide cada refinamiento sea coherente con la función objetivo del
co-agrupamiento, que distinga la heterogeneidad real del ruido y que preserve
la interpretabilidad del resultado, de modo que los patrones locales
detectados puedan leerse directamente como grupos de regiones y de horarios.

\subsection{Objetivos Específicos}
\label{subsec:objetivos-especificos}

\begin{enumerate}
    \item Formular el refinamiento como un árbol de bloques anidados que
          parte de la partición base de \ac{BBACI}, y definir un criterio
          de heterogeneidad por bloque coherente con la función objetivo 
          de \ac{ST-CC}.

    \item Diseñar la regla de decisión y de parada del refinamiento,
          combinando una prueba global por bloque, la selección del mejor
          corte en filas, en columnas o en ambas a la vez, pruebas calibradas
          con la variabilidad entre días de los datos, el control de
          comparaciones múltiples a lo largo del árbol y un umbral mínimo de
          tamaño de efecto.

    \item Implementar el método propuesto y evaluar su calibración y
          potencia sobre datos sintéticos con patrones locales conocidos y
          sobre datos nulos que reproducen la sobredispersión de los datos
          reales.

    \item Aplicar el método a los datos de taxis de Manhattan y compararlo
          con \ac{ST-CC} y con alternativas de rejilla global en términos
          de error de reconstrucción por número de parámetros, estabilidad
          e interpretabilidad de los patrones locales detectados.
\end{enumerate}

% ---------------------------------------------------------------------
\section{Organización de la tesis}
\label{sec:organizacion}

El restante de este documento de tesis está organizado de la siguiente
forma. Los capítulos siguen el orden natural del trabajo, que va de los
fundamentos teóricos a la propuesta, de la propuesta a su evaluación
experimental y, finalmente, a las conclusiones. Cada capítulo se resume a
continuación.

\begin{itemize}
    \item El Capítulo~\ref{chap:marco} presenta el marco teórico y el estado
          del arte, e incluye el co-agrupamiento basado en divergencias de
          Bregman y su lectura como modelo de Poisson, el marco \ac{ST-CC}, las
          pruebas de hipótesis para tablas de conteo, el agrupamiento guiado
          por pruebas de hipótesis (SigTree) y la revisión de los trabajos
          relacionados publicados desde 2020.
    \item El Capítulo~\ref{chap:propuesta} presenta la propuesta, que abarca la
          formulación del refinamiento jerárquico, el criterio estadístico y
          su calibración con la variabilidad entre días, la regla de parada y
          el algoritmo.
    \item El Capítulo~\ref{chap:resultados} expone los experimentos con
          datos sintéticos y reales, junto con el análisis de los
          resultados.
    \item El Capítulo~\ref{chap:conclusiones} resume las conclusiones, los
          problemas encontrados, las recomendaciones y los trabajos futuros.
\end{itemize}

% ---------------------------------------------------------------------
\section{Cronograma}
\label{sec:cronograma}

La Tabla~\ref{tab:cronograma} presenta la planificación de las actividades
de la tesis, organizadas en diez meses y vinculadas a los objetivos
específicos (OE) definidos en la Sección~\ref{subsec:objetivos-especificos}. Las
primeras actividades corresponden a la revisión de la literatura y a la
reproducción de \ac{ST-CC}, que sirven de base para el resto del trabajo. Luego
siguen el diseño y la implementación del método, su evaluación con datos
sintéticos y reales, y por último la redacción final y la revisión con el
asesor.

\begin{table}[htbp]
\centering
\small
\setlength{\tabcolsep}{3pt}  % requiere \usepackage{array} en el preámbulo
\begin{tabular}{>{\raggedright\arraybackslash}p{5.4cm} c *{10}{c}}
\hline
\textbf{Actividad} & \textbf{OE} &
\textbf{M1} & \textbf{M2} & \textbf{M3} & \textbf{M4} & \textbf{M5} &
\textbf{M6} & \textbf{M7} & \textbf{M8} & \textbf{M9} & \textbf{M10} \\
\hline
Revisión de la literatura y redacción del marco teórico
    & -- & $\bullet$ & $\bullet$ & & & & & & & & \\
Reproducción de \ac{ST-CC} sobre los datos de taxis de Manhattan
    & 4 & & $\bullet$ & $\bullet$ & & & & & & & \\
Formulación del árbol de bloques y del criterio de heterogeneidad
    & 1 & & $\bullet$ & $\bullet$ & & & & & & & \\
Diseño de la regla de decisión, parada y control de multiplicidad
    & 2 & & & $\bullet$ & $\bullet$ & & & & & & \\
Implementación del método propuesto
    & 3 & & & & $\bullet$ & $\bullet$ & $\bullet$ & & & & \\
Generación de datos sintéticos y experimentos de calibración y potencia
    & 3 & & & & & $\bullet$ & $\bullet$ & & & & \\
Experimentos con datos reales y comparación con métodos base
    & 4 & & & & & & $\bullet$ & $\bullet$ & & & \\
Análisis de estabilidad y visualización de resultados
    & 4 & & & & & & & $\bullet$ & $\bullet$ & & \\
Redacción de propuesta, resultados y conclusiones
    & -- & & & & & & $\bullet$ & $\bullet$ & $\bullet$ & $\bullet$ & \\
Revisión con el asesor y sustentación
    & -- & & & & & & & & & $\bullet$ & $\bullet$ \\
\hline
\end{tabular}
\caption{Cronograma de actividades de la tesis.}
\label{tab:cronograma}
\end{table}